\chapter{Bevezetés} % Introduction
\label{ch:intro}

Szakdolgozatom célja egy háromdimenziós, belső nézetes labirintusjáték megvalósítása C++ programozási nyelven. A megjelenítéshez OpenGL-t fogok használni segítségül. Az OpenGL egy olyan többnyelvű, multiplatformos alkalmazásprogramozási felület (API), amellyel közvetlenül, parancsokkal vezérelhetjük a grafikus processzort (azaz a GPU-t), ezáltal megvalósítható a 2D-s, vagy 3D-s alakzatok képernyőre rajzolása. Egy OpenGL kontextust többféleképpen is lehet implementálni, különböző programkönyvtárak segítségével. Ilyenek például a FreeGLUT, GLFW, SDL, vagy SFML.

Mivel a szakdolgozat egy játék megírására specializálódott, ezért ehhez az SDL (azaz Simple DirectMedia Layer) könyvtárat \cite{libsdl} fogom alkalmazni. Az SDL lényegében egy C nyelven alapuló multimédiás könyvtár, amellyel nem csak egy OpenGL kontextust lehet létrehozni, hanem lehetséges vele a bemeneti eszközök kezelése, hangok lejátszása, vagy a hálózatkezelés, a többjátékos (multiplayer) játékok elkészítéséhez.

A játékmotor az objektumorientált programozás (OOP) paradigmáját követi, így a program kódja az objektumok révén könnyen átlátható, valamint jól karbantartható és bővíthető. A főbb komponensek (például pályák, entitások, poligonhálók) különböző osztályokba szerveződnek. Az osztályokról részletesebben szó fog esni a fejlesztői dokumentációban. Az osztályok absztrakt reprezentációja miatt többféleképpen implementálható az egyes komponensek működése, így akár teljesen egyedi játékokat is tudunk készíteni. Egy labirintusjátékot fogok benne elkészíteni, amelyben a pályák gráfelméleti algoritmussal automatikusan generáltak, illetve egy beépített pályaszerkesztővel lehetőség van saját, egyedi pályák készítésére is. Egy 3D labirintusjáték szerintem nagyon jól szemlélteti a játékmotor adta lehetőségeket. A labirintusok tulajdonképpen pályák, amelyek entitásokból (játékosokból, falakból stb.) állnak és az entitásokhoz háromdimenziós modellek tartoznak, ezáltal a háromdimenziós térben tudjuk reprezentálni az elhelyezkedésüket és a megjelenésüket. 

A többjátékos mód megvalósításával bemutatom a TCP alapú kliens-szerver kommunikációt az \verb|SDL_net| hálózatkezelési könyvtár segítségével. A hálózati kommunikáció a kliensek (játékosok), valamint egy többszálú (multithreaded) szerverprogram között zajlik. Ezen játékmód kialakítása a játékélményhez is nagyban hozzájárul, ugyanis ekkor másokkal is játszhatunk egy játékot valós időben, ugyanazon a pályán.